\chapter{EJECUCIÓN Y VALIDACIÓN DEL SISTEMA}
\label{chap:ejecucion}

Este capítulo describe cómo poner en marcha el pipeline y cómo reproducir las mediciones presentadas en el capítulo anterior. Se ofrecen dos vías de arranque —una automática, pensada para la demostración, y una manual paso a paso— y se documentan las herramientas con las que se obtuvieron los resultados de los objetivos.

\section{Requisitos previos}
\label{sec:ejec-requisitos}

El sistema se ejecuta sobre Linux (o WSL2) y necesita:

\begin{itemize}
  \item \textbf{Docker} y \textbf{Docker Compose}, para el stack de servicios (Mosquitto, Kafka, Apicurio, TimescaleDB, PostgreSQL, Grafana y el bridge).
  \item \textbf{Java 21 LTS}, gestionado por SDKMAN a partir del fichero \texttt{.sdkmanrc} del repositorio. Lo requiere PySpark.
  \item \textbf{Python 3} con un entorno virtual (\texttt{venv}).
\end{itemize}

El script \texttt{bash setup.sh}, en la raíz del repositorio, comprueba la presencia de Docker y Java (guía la instalación, no la realiza), crea el entorno virtual, instala las dependencias de \texttt{requirements.txt} y prepara el conjunto de datos ASHRAE si encuentra credenciales de Kaggle. El entorno virtual debe estar activado antes de lanzar los procesos que corren en el host: el job de Spark y el simulador.

\section{Modo automático (demostración)}
\label{sec:ejec-demo}

El script \texttt{pipeline/tools/demo.py} orquesta todo el arranque en un solo comando: levanta el stack, deja las bases de datos limpias, arranca el bridge y el job de Spark —esperando a que las consultas de \textit{streaming} estén activas— y lanza el simulador reproduciendo los datos más recientes del histórico.

{\footnotesize
\begin{verbatim}
python pipeline/tools/demo.py            # 6 semanas de datos, arranque completo
python pipeline/tools/demo.py --semanas 10
python pipeline/tools/demo.py --stop     # cierra los procesos y baja el stack
\end{verbatim}
}

La opción \texttt{-{}-semanas N} reproduce la cola de las últimas \texttt{N} semanas del histórico (entre 1 y 52; 6 por defecto). Como el Parquet de datos ya viene fechado en el presente, esa cola termina en fecha reciente. Mientras el simulador publica, el dashboard de Grafana (\texttt{http://localhost:3000}) se rellena de izquierda a derecha; en los dashboards operacionales conviene fijar el rango temporal a \texttt{now-Nw} para ver el bloque completo.

Esta vía sirve para \textbf{demostrar} el sistema en funcionamiento, no para medir: las cifras de los KPIs se toman con el procedimiento manual de la sección siguiente.

\section{Modo manual}
\label{sec:ejec-manual}

El arranque manual sigue cinco pasos, en este orden:

\begin{enumerate}
  \item \textbf{Estado limpio.} Borra los \textit{checkpoints} de Spark, recrea los tópicos de Kafka y vacía las tablas de los sumideros, de modo que la medición no se mezcle con los restos de una ejecución anterior.
  {\footnotesize
  \begin{verbatim}
python pipeline/tools/reset_state.py --yes
  \end{verbatim}
  }

  \item \textbf{Levantar el stack.} Arranca los siete servicios. El registro del esquema (\texttt{register-schema}, un contenedor de un solo uso) y el bridge están encadenados por \texttt{depends\_on}: el bridge no arranca hasta que el esquema Avro está registrado.
  {\footnotesize
  \begin{verbatim}
docker compose -f pipeline/docker-compose.yml up -d
  \end{verbatim}
  }

  \item \textbf{Arrancar el job de Spark.} El primer arranque descarga los conectores de Maven y tarda más de lo habitual; hay que esperar a que el job informe de que las consultas de \textit{streaming} están en marcha antes de generar carga.
  {\footnotesize
  \begin{verbatim}
python pipeline/spark/stream_processing.py --trigger "1 second"
  \end{verbatim}
  }

  \item \textbf{Generar la carga.} El simulador reproduce la telemetría de los 652 sensores. El factor \texttt{-{}-acelerar} comprime el reloj del histórico.
  {\footnotesize
  \begin{verbatim}
python pipeline/simulator/mqtt_simulator.py --acelerar 2000
  \end{verbatim}
  }

  \item \textbf{Emitir el cuadro de KPIs.} Consulta las fuentes donde el pipeline deja constancia de su actividad y escribe el informe en \texttt{pipeline/logs/informe\_kpi.md}.
  {\footnotesize
  \begin{verbatim}
python pipeline/tools/kpi_report.py
  \end{verbatim}
  }
\end{enumerate}

El orden importa: el job de Spark debe estar procesando \textbf{antes} de que el simulador publique. Si se invierte, los eventos esperan en Kafka a que aparezca quien los consuma y esa espera se contabiliza como latencia de ingesta, falseando la medición.

\section{Herramientas de validación}
\label{sec:ejec-herramientas}

Las mediciones de los objetivos se obtuvieron con tres scripts que no forman parte del pipeline —no procesan telemetría, la miden o preparan el sistema para medirla— y que residen en \texttt{pipeline/tools/}.

\begin{table}[H]
  \centering
  \begin{tabular}{|p{3.2cm}|p{7.3cm}|p{3.3cm}|}
    \hline
    \textbf{Herramienta} & \textbf{Función} & \textbf{Salida} \\ \hline
    \texttt{reset\_state.py} & Deja el sistema en estado limpio: borra los \textit{checkpoints} de Spark, recrea los tópicos de Kafka y vacía las tablas de los sumideros. & --- \\ \hline
    \texttt{kpi\_report.py} & Interroga PostgreSQL, TimescaleDB, Apicurio y Grafana y emite el cuadro de los KPIs 1 a 5 en Markdown. & \texttt{pipeline/logs/} \texttt{informe\_kpi.md} \\ \hline
    \texttt{failover\_test.py} & Lanza el simulador, provoca el fallo de un servicio y mide el tiempo de recuperación y la tasa de pérdida (Objetivo 5). & \texttt{pipeline/logs/} \texttt{ultimo\_failover.json} \\ \hline
  \end{tabular}
  \caption{Herramientas de validación del sistema.}
  \label{tab:herramientas-validacion}
\end{table}

\texttt{kpi\_report.py} no mide una prueba concreta, sino el estado en que encuentra el sistema. Para que sus cifras sean fiables debe ejecutarse siguiendo el orden estricto de la sección~\ref{sec:ejec-manual} (estado limpio, stack, Spark, carga, informe); la alternativa es apoyarse en \texttt{demo.py} (sección~\ref{sec:ejec-demo}), que realiza esa misma secuencia de puesta en marcha de forma automática.

\texttt{failover\_test.py} necesita que se haya ejecutado \texttt{reset\_state.py} antes de arrancar el pipeline: el simulador que lanza internamente parte de un estado limpio. La prueba comprueba que el bridge y el job de Spark estén en marcha y se detiene con un mensaje si falta alguno. El servicio a interrumpir se indica con \texttt{-{}-target} (Mosquitto, Kafka, TimescaleDB, PostgreSQL o el bridge).

{\footnotesize
\begin{verbatim}
python pipeline/tools/failover_test.py --target postgres
\end{verbatim}
}
